\chapter{La fisica del movimento}
\label{cap:fisica_movimento}

Il capitolo precedente ha lasciato una domanda aperta: quando un modello basato sul solo sforzo energetico è davvero applicabile a un percorso antico? Per rispondere, occorre prima capire cosa dice davvero la scienza del movimento su pendenza, e se le leggi che governano lo sforzo di un essere umano siano le stesse che governano lo sforzo di un animale selvatico.

\section{Una funzione nata per gli escursionisti}

Negli studi di geografia ed escursionismo esiste da tempo una formula, oggi diffusamente utilizzata anche in archeologia, che stima la velocità di cammino di una persona in funzione della pendenza del terreno: più il terreno sale o scende ripidamente, più la velocità cala, con un andamento che non è simmetrico tra salita e discesa (scendere è meno faticoso che salire, fino a un certo punto, oltre il quale anche la discesa diventa lenta e difficoltosa). Questa formula è stata sviluppata originariamente per descrivere il comportamento di camminatori moderni su terreno generico, e la sua semplicità l'ha resa uno strumento popolare per stimare, da una mappa altimetrica, quanto tempo servirebbe a percorrere un dato tragitto a piedi.

Il problema, quando si vuole applicare questa formula a un contesto preistorico, è capire quanto sia legittimo trasportarla di peso da un escursionista moderno, con scarpe adeguate e senza zavorra particolare, a un cacciatore di diecimila anni fa, o addirittura a un animale selvatico. È una domanda che merita una risposta basata su dati, non su intuizione.

\section{Un modello più recente, e più preciso}

La scienza del movimento ha affrontato esplicitamente questa domanda, e lo ha fatto in modo particolarmente rigoroso con un lavoro pubblicato nel 2016 dal biologo Herman Pontzer su \textit{Biology Letters} \citep{pontzer2016}. Pontzer propone un modello, oggi noto come modello ARC (dalle iniziali di \textit{Activation-Relaxation and Cross-bridge cycling}, i due processi della fisiologia muscolare su cui il modello si fonda), che descrive il costo energetico della locomozione terrestre in funzione della pendenza, derivandolo dai meccanismi cellulari della contrazione muscolare invece che da una semplice osservazione empirica del comportamento.

Il punto di forza di questo modello, per il problema di questo libro, è l'ampiezza della sua validazione: Pontzer lo ha testato su un intervallo di masse corporee che va da 0,78 grammi a 431 chilogrammi, e su un intervallo di pendenze che va da meno ventiquattro gradi (una discesa moderata) a novanta gradi (l'arrampicata verticale), coprendo un ventaglio di specie che comprende insetti, piccoli mammiferi e grandi ungulati. Il modello descrive con successo il comportamento osservato sia all'interno di ciascuna specie sia nel confronto tra specie diverse, per l'intero ventaglio di pendenze testate, incluso il caso limite dell'arrampicata verticale.

Un lavoro successivo, che ha sviluppato uno strumento software per la modellazione del paesaggio energetico (\termine{enerscape}), ha applicato il modello ARC a scenari con specie di taglia variabile, incluse applicazioni a grandi mammiferi, confermandone la coerenza su un ampio spettro di scale \citep{berti2022}. Il modello non è stato tarato su misura per un solo gruppo di taglia corporea, e sembra generalizzare bene anche fuori dal range su cui è stato costruito.

\section{Cosa cambia tra specie, e cosa resta uguale}

Un secondo lavoro di sintesi, pubblicato lo stesso anno da Lewis Halsey su \textit{Journal of Experimental Biology}, offre un quadro complementare: per la maggior parte delle specie, il tasso a cui l'energia viene consumata durante la locomozione è, con buona approssimazione, correlato linearmente alla velocità di movimento, almeno su terreno pianeggiante \citep{halsey2016}. Anche questo lavoro conferma che la relazione tra costo energetico e velocità (o pendenza) segue pattern simili tra specie molto diverse, pur variando nei parametri specifici.

Ciò che cambia da specie a specie, secondo entrambi questi lavori, non è la forma della relazione, ma i suoi parametri specifici, legati soprattutto alla massa corporea e alla lunghezza degli arti: gli animali di piccola taglia, con arti corti e passi più rapidi, sostengono un costo energetico proporzionalmente maggiore rispetto agli animali di taglia più grande, anche a causa dei processi di attivazione e rilassamento muscolare che si ripetono più frequentemente. Tra specie di massa corporea simile, però, la differenza nei parametri della curva è ridotta.

Per il problema di questo libro, questo è un dato rassicurante. Gli ungulati alpini di interesse discussi nel quinto capitolo (stambecco, camoscio, cervo) hanno una massa corporea nello stesso ordine di grandezza di un essere umano adulto, tra i quaranta e i novanta chilogrammi circa a seconda della specie e dell'individuo, ben all'interno del range di 0,78 grammi - 431 chilogrammi su cui il modello ARC è stato validato. Non c'è quindi, in linea di principio, una ragione fisiologica forte per aspettarsi un comportamento radicalmente diverso da quello descritto dal modello: la fisica di base, il lavoro necessario per sollevare il proprio peso contro la gravità lungo un pendio, non cambia sostanzialmente tra un cacciatore e la sua preda, a parità di taglia corporea.

\section{Dove la fisica basta, e dove non basta più}

A questo punto è possibile chiudere il cerchio aperto nel capitolo precedente. La fisica del movimento su pendenza, nella sua formulazione più aggiornata e meglio validata, descrive bene il costo energetico degli spostamenti di un animale selvatico: il costo della pendenza è uno dei principali vincoli che condizionano la scelta del percorso, insieme a habitat, neve, cibo, riproduzione e rischio di predazione. Non ci sono, per un camoscio o per uno stambecco, ragioni simboliche o amministrative che possano spiegare uno scarto sistematico tra il percorso atteso e quello effettivo. In questo caso, un modello di costo energetico come quello di Pontzer è uno strumento legittimo, con solide basi scientifiche e un'ampiezza di validazione che copre esattamente il range di taglia corporea che interessa questo libro.

La stessa fisica, applicata a un percorso umano pianificato, ingegnerizzato, o carico di significati che vanno oltre il puro spostamento fisico (come il Qhapaq Ñan discusso nel capitolo precedente, o persino i sentieri pedonali gallesi studiati da Llobera e Sluckin), smette di essere sufficiente da sola: descrive ancora, correttamente, quanto sarebbe costato energeticamente percorrere una data via, ma non può spiegare perché gli uomini abbiano scelto, in tanti casi concreti, di percorrerne un'altra.

Questa distinzione ha una conseguenza diretta e importante per il problema di questo libro. Se l'obiettivo è ricostruire dove si muovevano gli animali cacciati (stambecco, camoscio, cervo, e nelle fasi più antiche uro, bisonte, cavallo), un modello di costo energetico basato sulla pendenza, e in particolare il modello ARC nella sua formulazione più recente, è uno strumento ragionevole. Se l'obiettivo, invece, è ricostruire dove si muovevano i cacciatori umani, la questione si complica: i cacciatori inseguivano le prede, quindi in parte ne condividevano i tragitti, ma portavano con sé anche altre necessità (accesso all'acqua per sé e per il gruppo, riparo dalle intemperie, prossimità a un accampamento condiviso con il resto della famiglia allargata) che un modello di sola pendenza non cattura.

\section{Una tentazione che merita di essere raccontata}

A questo punto della ricerca che ha portato a scrivere questo libro, la tentazione più naturale era costruire un modello che seguisse esattamente questa logica: calcolare, da una mappa del terreno, i percorsi a minor costo energetico per un animale di taglia media, individuare dove questi percorsi convergono verso punti di passaggio obbligato, e proporre quei punti come i luoghi più probabili per un accampamento di caccia preistorico.

È un'idea che sembra, a prima vista, tanto elegante quanto ben fondata scientificamente, alla luce di quanto appena discusso sulla fisiologia del movimento. Il prossimo capitolo racconta perché, messa alla prova dei dati reali sul comportamento degli animali alpini già presentati nel quinto capitolo, questa idea si sia rivelata meno solida di quanto sembrasse, e cosa quell'errore abbia insegnato per il seguito del ragionamento.
